% ============================================================================
% GENERATION AS A BOUNDARY VALUE PROBLEM
% ============================================================================
\subsection{The Autoregressive IVP Limitation}
\label{sec:ivp_limitation}

Contemporary language models formulate generation as an autoregressive
Initial Value Problem (IVP): given a starting context, continuously predict
the next token.  Formally, given prefix $\mathbf{x}_{0:t}$, the IVP
computes $p(x_{t+1} \mid \mathbf{x}_{0:t})$ and samples, yielding a
Markov chain with strictly left-to-right information flow.  This creates
three structural limitations:

\begin{enumerate}[nosep]
  \item \textbf{Serial depth bottleneck:}  Each step depends on the
        materialization of all preceding outputs, forcing an $O(N)$
        critical-path depth that hardware parallelism cannot bypass.
  \item \textbf{Irrevocable commitment:}  Tokens cannot be retracted once
        emitted; local errors propagate as distributional drift---a single
        decoding error permanently alters the generation trajectory.
  \item \textbf{Static memory:}  Knowledge is compressed into weight
        matrices $\mathbf{W}$ requiring gradient descent to update,
        causing catastrophic forgetting.
\end{enumerate}

\subsection{BVP Formulation via Topological Vacuum}
\label{sec:topological_vacuum}

We frame generation as a \emph{Boundary Value Problem solved through
discrete geometric resonance}.  Given boundary constraints encoded as
a chord $\mathbf{e}$ (the existing context) and a target distribution
$\mathbf{t}$ (either the context itself or an ExpectedReality manifold),
the system computes the Topological Vacuum:

\begin{equation}
  \mathbf{h} = H(\mathbf{t}, \mathbf{e}) = \mathbf{t} \land \lnot\,\mathbf{e}
  \label{eq:chord_hole_bvp}
\end{equation}

This isolates the exact prime oscillators \emph{required but missing} from
the context.  The active bits in $\mathbf{h}$ represent structural
constraints that must be satisfied by the continuation; bits not in
$\mathbf{h}$ are already satisfied or irrelevant.

\subsection{Resonant Fill: The FillScore}
\label{sec:resonant_fill}

Rather than computing transition probabilities, the inference engine
searches the entire PrimeField for the manifold block whose topology most
perfectly \emph{fills} the vacuum.  Given vacuum $\mathbf{h}$ and
candidate chord $\mathbf{d}$:

\begin{equation}
  \text{FillScore}(\mathbf{h}, \mathbf{d}) =
    \frac{\texttt{popcount}(\mathbf{h} \land \mathbf{d})}
         {\texttt{popcount}(\mathbf{h})
          \;\cdot\;
          \bigl(1 + \texttt{popcount}(\mathbf{d} \land \lnot\,\mathbf{h})\bigr)}
  \label{eq:fill_score}
\end{equation}

\noindent The three terms have precise physical interpretations:

\begin{itemize}[nosep]
  \item \textbf{Numerator} $\texttt{popcount}(\mathbf{h} \land \mathbf{d})$:
        The number of vacuum bits filled by the candidate---\emph{constructive
        interference} between the structural need and the candidate's content.
  \item \textbf{Denominator term 1} $\texttt{popcount}(\mathbf{h})$:
        The total vacuum size, normalizing the score to $[0, 1]$.
  \item \textbf{Denominator term 2}
        $1 + \texttt{popcount}(\mathbf{d} \land \lnot\,\mathbf{h})$:
        A penalty for \emph{structural noise}---candidate oscillators not
        demanded by the vacuum.  The $+1$ prevents division by zero and
        ensures that a candidate with zero noise achieves the maximum score.
\end{itemize}

\paragraph{Properties of FillScore.}
\begin{enumerate}[nosep]
  \item \textbf{Bounded:} $\text{FillScore} \in [0, 1]$.
        Achieved at 1 iff $\mathbf{d}$ fills every vacuum bit and introduces
        no extra bits ($\mathbf{h} \land \mathbf{d} = \mathbf{h}$ and
        $\mathbf{d} \land \lnot\mathbf{h} = \mathbf{0}$).
  \item \textbf{Monotone in fill:}  For fixed $\mathbf{h}$, adding a
        vacuum-aligned bit to $\mathbf{d}$ strictly increases the score.
  \item \textbf{Penalizing noise:}  For fixed fill, adding a non-vacuum bit
        to $\mathbf{d}$ strictly decreases the score via the denominator
        penalty.
  \item \textbf{$O(D/w)$ computation:}  All three popcount terms are
        computed in a single pass over the $D/64 = 8$ blocks.
\end{enumerate}

\subsection{Transitive Resonance: Analogical Transfer}
\label{sec:transitive_resonance}

The chord algebra supports structural analogy without neural weights.
Given three facts encoded as chords $F_1, F_2, F_3$, the system
extracts shared factors and forges a hypothesis:

\begin{align}
  B &= F_1 \land F_2 && \text{(Shared context: ``wormhole'')} \label{eq:tr_b} \\
  C &= F_2 \land F_3 && \text{(Shared subject)} \label{eq:tr_c} \\
  A &= F_1 \land \lnot\,B && \text{(Unique to } F_1 \text{)} \label{eq:tr_a} \\
  D &= F_3 \land \lnot\,C && \text{(Unique to } F_3 \text{)} \label{eq:tr_d} \\
  H &= A \lor D && \text{(Forged hypothesis)} \label{eq:tr_h}
\end{align}

\noindent This implements the analogy $A : B :: C : D$ using five bitwise
operations in $O(D/w)$ time, enabling zero-shot relational transfer.

\subsection{The Cantilever Generation Loop}
\label{sec:cantilever}

Generation extends the original prompt base like a physical cantilever:

\begin{enumerate}[nosep]
  \item Build the context manifold $\mathbf{e}_0$ from the prompt chords.
  \item Compute the structural vacuum $\mathbf{h}_t = H(\mathbf{t},
        \mathbf{e}_t)$.
  \item Execute BestFill across all 5 Fibonacci planes simultaneously,
        returning $\hat{\mathbf{d}} = \arg\max_{\mathbf{d} \in \text{PF}}
        \text{FillScore}(\mathbf{h}_t, \mathbf{d})$.
  \item Check EigenMode intent alignment:
        $\Delta\theta(\hat{\mathbf{d}}) < \delta_{\text{intent}}$.
  \item Integrate: $\mathbf{e}_{t+1} = \mathbf{e}_t \lor
        \hat{\mathbf{d}}$.
  \item Emit the byte-level readout of $\hat{\mathbf{d}}$.
  \item Repeat until geometric closure is detected (\Cref{sec:geometric_closure})
        or $t > t_{\max}$.
\end{enumerate}

The growing structure becomes its own foundation for subsequent searches,
forming a positive feedback loop where each successful fill sharpens the
context for the next.
